\section{Introduction}
\label{sec:introduction}

% TODO(figure): add the teaser \ref{fig:teaser} (escalation vs. continuous attack
% under a mean-pooling detector). TODO(results): fill the [XX] placeholders.

LLM-based multi-agent systems (MAS) now solve tasks that exceed the reach of any
single agent, ranging from collaborative reasoning and debate to tool use and
autonomous decision-making~\citep{guo2024survey,li2023camel,zhuge2024gptswarm}.
Their strength comes from inter-agent communication, but this same communication
is a new attack surface: a few compromised agents can inject malicious content
that propagates across the network and degrades the whole system rather than a
single node~\citep{yu2024netsafe}. Such compromises stem from familiar vectors
---direct and indirect prompt injection~\citep{perez2022ignore,greshake2023indirect},
tool attacks~\citep{zhan2024injecagent}, and memory poisoning~\citep{chen2024agentpoison}---
yet in a MAS their damage is amplified by propagation instead of being contained
at the point of entry.

To defend MAS, recent work models the system as a graph and detects malicious
agents from their communication. The state-of-the-art defense,
G-Safeguard~\citep{wang2025gsafeguard}, constructs a multi-agent \emph{utterance
graph} whose nodes are agents and whose edges carry the messages exchanged
between them, trains a graph neural network (GNN) to flag high-risk agents, and
prunes their edges to halt contamination. Because the GNN is inductive, a single
detector transfers across LLM backbones, topologies, and system scales without
retraining. This detect-and-remediate paradigm is currently the most general and
effective line of defense for MAS.

However, such detectors collapse the temporal dimension of multi-turn dialogue
before they reason over the graph. Each edge accumulates an agent's utterances
across several rounds, but the detector compresses this sequence into one
embedding by mean-pooling---or by keeping only the last turn---before message
passing. Mean-pooling discards \emph{how} an agent's behavior evolves over the
conversation, and we identify this as a critical blind spot: detection then
depends on an agent's average content rather than on its trajectory. An attacker
can exploit this directly. If it behaves benignly in the early turns and escalates
its malicious intent only gradually, the benign early turns dominate the average
and pull the attacker's mean embedding toward that of honest agents---saturating
the very feature on which detection relies, so the attacker evades the detector
while still steering the system.

We make two contributions that turn this blind spot into both a threat and a
remedy. First, we introduce the \emph{escalation attack}, a stealthy strategy in
which a compromised agent starts indistinguishable from honest agents and ramps
up its influence over successive turns (Figure~\ref{fig:teaser}). By saturating
the mean-pooled feature with benign early turns, the escalation attack sharply
lowers the detection rate of existing defenses while remaining effective at
compromising the MAS. Second, we propose \textbf{T-GAT}, a temporal-aware
detector that replaces mean-pooling with a temporal edge encoder. Rather than
averaging an agent's utterances, T-GAT extracts their temporal structure through
self-attention over turns, a temporal convolution that captures local shifts, and
explicit cross-turn statistics---variance and step-wise change---that directly
expose escalation; a learned gate then fuses these temporal cues with message
content before graph message passing.

The key insight behind T-GAT is that the discriminative signal is \emph{how} an
agent changes, not only \emph{what} it says on average. An honest agent is
consistent across turns, so its cross-turn variance is near zero; an escalating
attacker is, by construction, inconsistent, so the statistics a mean obscures
become its signature. Because T-GAT preserves these dynamics instead of discarding
them, it is designed to detect escalation attacks that defeat mean-pooling while
remaining effective against continuous attacks, whose strong per-turn signals
existing detectors already capture. T-GAT keeps the inductive, drop-in design of
graph-based detectors, so it inherits their transferability across backbones and
topologies.

We evaluate across diverse agent-attack scenarios, multiple LLM backbones, and
graph topologies. The escalation attack reduces the detection accuracy of
G-Safeguard by \textbf{[XX]} and raises the attack success rate by \textbf{[XX]},
while T-GAT recovers detection to \textbf{[XX]} and lowers the attack success rate
by \textbf{[XX]} under both continuous and escalation attacks.
% >>NEEDS<< finalize the exact attack scenarios / datasets evaluated and cite them,
% e.g., tool attack (InjecAgent), prompt injection (CSQA/MMLU/GSM8K), memory
% poisoning (PoisonRAG). Keep this consistent with the Abstract and Experiments.

\noindent\textbf{Contributions.} We summarize our contributions as follows:
\begin{itemize}
    \item We expose a temporal blind spot in state-of-the-art graph-based MAS
    defenses: by mean-pooling multi-turn dialogue, attacker detection depends on
    average content and ignores behavioral dynamics.
    \item We introduce the \emph{escalation attack}, which exploits this blind
    spot by saturating the mean-pooled feature with benign early turns, evading
    detection while still compromising the MAS.
    \item We propose \textbf{T-GAT}, a temporal-aware detector that models
    cross-turn dynamics---self-attention, temporal convolution, and
    variance/step-change statistics---and defends against both escalation and
    continuous attacks while preserving inductive transferability.
    \item Experiments across multiple LLM backbones and graph topologies show that
    [\textbf{summary of results once measured}].
\end{itemize}
